% =====================================================================
% agora-evidence-chain-paper.tex
% Wissenschaftliches Forschungs-Paper aus dem Agora Evidence-Chain-Audit
% Stand: 2026-08-09 (Repo-HEAD 7e42ae34, Branch feat/1152-provenance-retrieval)
% Engine: LuaLaTeX / BibLaTeX+Biber / KOMA-Script scrreprt
% =====================================================================

\documentclass[11pt,a4paper,oneside,open=any,parskip=half-]{scrreprt}

% --- Sprache ---------------------------------------------------------
\usepackage{polyglossia}
\setmainlanguage{german}
\setotherlanguage{english}

% --- Schriften (LuaLaTeX + fontspec) ---------------------------------
\usepackage{fontspec}
\setmainfont{Libertinus Serif}[
  Numbers={OldStyle,Proportional},
  Ligatures=TeX
]
\setsansfont{Libertinus Sans}[
  Scale=MatchUppercase
]
\setmonofont{Libertinus Mono}

% --- Symbole (Kreuz, Haken) ------------------------------------------
\usepackage{pifont}
\newcommand{\lossmark}{\texorpdfstring{\ding{55}}{\ding{55}}}
\newcommand{\okmark}{\texorpdfstring{\ding{51}}{\ding{51}}}
\newcommand{\partialmark}{\texorpdfstring{\ding{74}}{\ding{74}}}
% Subscript-Helper fuer Tabellen: \lossmarkX{1} -> ✗₁ (Pifont-Font erzwungen)
\newcommand{\lossmarkX}[1]{\texorpdfstring{\ding{55}\textsubscript{#1}}{X-#1}}

% --- Mikrotypografie --------------------------------------------------
\usepackage{microtype}
\usepackage{csquotes}

% --- Mathe -----------------------------------------------------------
\usepackage{amsmath,amssymb}

% --- Tabellen --------------------------------------------------------
\usepackage{booktabs}
\usepackage{tabularray}

% --- Grafik / TikZ ---------------------------------------------------
\usepackage{graphicx}
\usepackage{tikz}
\usetikzlibrary{positioning,arrows.meta,shapes.geometric,fit,calc,backgrounds,shadows.blur,matrix}

% --- Hyperref (zuletzt) ----------------------------------------------
\usepackage[
  bookmarks=true,
  bookmarksopen=true,
  bookmarksnumbered=true,
  hidelinks,
  pdftitle={Agora Evidence Chain -- Ein Code-Audit zu Provenance und Anti-Self-Confirmation in einer Multi-Agent-Simulationsplattform},
  pdfauthor={Alexander Schneider},
  pdfsubject={Code-Audit, Provenance, Evidence-Gating, Multi-Agent-Simulation},
  pdfkeywords={Agora, Evidence-Gating, Provenance, Multi-Agent, Anti-Self-Confirmation, Code-Audit}
]{hyperref}

% --- Bibliographie ---------------------------------------------------
\usepackage[
  backend=biber,
  style=numeric-comp,
  sorting=none,
  citestyle=numeric-comp,
  maxnames=6,
  minnames=1,
  giveninits=true,
  isbn=false,
  doi=true,
  eprint=true,
  eprinttype=arxiv
]{biblatex}
\addbibresource{references.bib}

% --- Befehle ---------------------------------------------------------
\newcommand{\agora}{\textsc{Agora}}
\newcommand{\code}[1]{\texttt{#1}}
\newcommand{\audit}{\texttt{agora-evidence-chain-audit.md}}
\newcommand{\citesource}[1]{\hyperlink{src.#1}{[#1]}}
\newcommand{\src}[2]{\hypertarget{src.#1}{}\textbf{[#1]}\footnote{\raggedright #2}}

% Verbatim-Umgebung fuer Code-Snippets
\usepackage{fancyvrb}

% Schatten fuer TikZ-Boxen
\usetikzlibrary{shadows.blur}

% Quellen-Kurztabelle (Belege)
\newcommand{\beleg}[2]{\hyperlink{src.#1}{\textbf{[#1]}} & \small #2}

% --- Metadaten -------------------------------------------------------
\title{\agora{} Evidence Chain\\\large Ein Code-Audit zu Provenance und Anti-Self-Confirmation in einer Multi-Agent-Simulationsplattform}
\author{Alexander Schneider\\
\small \href{https://orcid.org/0009-0008-7552-1819}{ORCID 0009-0008-7552-1819}}
\date{2026-08-09 \\
\small Repository \code{arn0ld87/agora} -- \href{https://github.com/arn0ld87/agora}{github.com/arn0ld87/agora}, Branch \code{feat/1152-provenance-retrieval}, HEAD \code{7e42ae34} \\
\small Version 0.9.3 Stability Beta \\
\small Begleitendes Software-Artefakt: \href{https://zenodo.org/records/21855023}{Zenodo 21855023}}

\begin{document}

\maketitle
\thispagestyle{empty}

\begin{abstract}
\noindent
Diese Arbeit untersucht am Code-Audit eines realen Multi-Agent-Systems (\agora{}, Version 0.9.3 Stability Beta), welche Bestandteile der Evidence-Pipeline einen \emph{tats{\"a}chlich {\"u}berpr{\"u}fbaren} Erkenntnisgewinn gegen{\"u}ber einem einzelnen starken Large-Language-Model-Prompt erzeugen, und welche lediglich zus{\"a}tzliche Komplexit{\"a}t, Kosten oder scheinbare Evidenz. Methodisch verkn{\"u}pft der Beitrag eine statische Code-Audit-Pipeline (acht parallele Subagenten mit Code-Review-Graph-Unterst{\"u}tzung) mit einer normativen Bewertung anhand eines zehn-dimensionalen Scorerasters. Im Mittelpunkt steht die Rekonstruktion der Evidence-Chain vom Seed-Dokument bis zum Claim {\"u}ber 17 Stufen, einschließlich sieben code-verifizierter Provenance-Verlustpunkten. Die Ergebnisse zeigen, dass \agora{} als \emph{ehrlich konstruiertes, aber empirisch unbewiesenes} System zu lesen ist: Der Engineering-Score liegt bei 7,0/10, der Evidence-Research-Score bei 3,5/10; ein Referenzlauf lieferte null validierte Claims. Die Diskussion arbeitet f{\"u}nf Kernkontroversen heraus und priorisiert P0-Maßnahmen (Seed-Provenance, Interview-Binding, Konzepttrennung der Confidence-Scores), die den Mehrwert der Pipeline {\"u}berf{\"u}hrbar machen sollen. Der Beitrag ist sowohl Diagnose als auch Bauplan.
\end{abstract}

\begin{abstract}[english]
\noindent
This paper investigates, via a code audit of a real multi-agent system (\agora{}, v0.9.3 Stability Beta), which parts of the evidence pipeline yield a \emph{verifiable} epistemic gain over a single strong large language model prompt, and which add complexity, cost, or apparent evidence only. Methodologically, the work links a static code-audit pipeline (eight parallel sub-agents, supported by a code-review-graph) with a normative evaluation along a ten-dimensional scoring scheme. The core artefact is a reconstruction of the evidence chain from seed document to claim across seventeen stages, including seven code-verified provenance loss points. The findings read \agora{} as an \emph{honestly engineered but empirically unproven} system: engineering score 7.0/10, evidence-research score 3.5/10, one reference run producing zero validated claims. The discussion isolates five key controversies and prioritises P0 measures (seed-document provenance, interview-binding fix, separation of confidence concepts) that should make the pipeline's added value observable. The contribution is simultaneously diagnosis and blueprint.
\end{abstract}

\tableofcontents
\listoffigures
\listoftables

% =====================================================================
\part{Kontext und Forschungsfrage}
% =====================================================================

% =====================================================================
% Kapitel 1: Einleitung und Forschungsfrage
% =====================================================================

\chapter{Einleitung und Forschungsfrage}\label{ch:einleitung}

\section{Motivation}\label{sec:motivation}

Multi-Agent-Frameworks auf Basis großer Sprachmodelle werden zunehmend als Werkzeug für die Marktanalyse, Stakeholder-Simulation und Szenarienmodellierung eingesetzt~\cite{park2023generative,argyle2023silicon,horton2023homosilicus}. Die Versprechen sind hoch: synthetische Populationen sollen reale Zielgruppen approximieren, Multi-Agent-Debate soll die Factuality~\cite{du2023multiagent} erhöhen, und Selbstkonsistenz~\cite{wang2022selfconsistency} soll gegen Halluzination absichern. Die Literatur zeigt jedoch zugleich, dass Konsistenz nicht mit Korrektheit gleichgesetzt werden darf~\cite{consistency2024}, dass rekursives Training auf synthetischen Daten zu Model-Collapse führt~\cite{shumailov2024curse}, dass Inline-Zitationen häufig \enquote{cited but not verified} sind~\cite{citednotverified} und dass die Validität von LLM-Personas gegen reale Populationen eine offene empirische Frage bleibt~\cite{uas2025,analyticflexibility}.

Diese Arbeit setzt genau dort an: Sie untersucht am konkreten Code eines realen Multi-Agent-Systems (\agora{}), welche Bausteine der Evidence-Pipeline tatsächlich einen \emph{prüfbaren} Mehrwert liefern -- und welche lediglich zusätzliche Komplexität erzeugen, ohne den Anspruch einzulösen.

\section{Forschungsfrage}\label{sec:forschungsfrage}

Die zentrale Forschungsfrage lautet:

\begin{quote}
Welche Teile von \agora{} erzeugen einen tatsächlich überprüfbaren Erkenntnisgewinn gegenüber einem einzelnen starken LLM-Prompt -- und welche Teile erzeugen lediglich zusätzliche Komplexität, Kosten oder scheinbare Evidenz?
\end{quote}

Drei Unterfragen strukturieren die Antwort:

\begin{enumerate}
  \item \textbf{Provenance:} Lässt sich der Pfad vom Seed-Dokument bis zum ausgelieferten Claim vollständig rekonstruieren -- und an welchen Stellen geht die Herkunft verloren?
  \item \textbf{Anti-Self-Confirmation:} Welche Mechanismen verhindern, dass synthetische Populationen sich gegenseitig bestätigen, ohne dass diese Mechanismen kalibriert oder extern validiert wären?
  \item \textbf{Empirischer Mehrwert:} Lässt sich an einem realen Referenzlauf zeigen, dass das System über einen Single-Prompt hinaus zusätzliche, \emph{geprüfte} Aussagen liefert -- oder erstickt es an seiner eigenen Strenge?
\end{enumerate}

\section{Beitrag}\label{sec:beitrag}

Der Beitrag dieser Arbeit ist dreigliedrig:

\begin{itemize}
  \item \textbf{Methodisch} wird eine reproduzierbare Code-Audit-Pipeline vorgestellt, die parallele Subagenten mit Code-Review-Graph-Unterstützung und einer normativen Bewertung verknüpft. Die Pipeline ist nicht auf \agora{} beschränkt; sie ist auf jedes Multi-Agent-System übertragbar, dessen Code zugänglich ist.
  \item \textbf{Diagnostisch} liefert die Arbeit eine Rekonstruktion der Evidence-Chain vom Seed-Dokument bis zum Claim, mit sieben code-verifizierten Provenance-Verlustpunkten, einer 10-Dimensionen-Bewertung und einer priorisierten Roadmap.
  \item \textbf{Normativ} wird ein Antwortrahmen vorgeschlagen, mit dem sich die Frage \enquote{Liefert ein LLM-System einen messbaren Mehrwert?} überhaupt erst formulieren lässt -- inklusive der Trennung in \emph{engineering score} (Konstruktion) und \emph{evidence-research score} (Empirie).
\end{itemize}

\section{Gliederung}\label{sec:gliederung}

Die Arbeit gliedert sich in sieben Kapitel. Kapitel~\ref{ch:related} arbeitet den Stand der Forschung anhand von fünf Themenclustern auf und benennt die Forschungslücke, in die der Beitrag fällt. Kapitel~\ref{ch:system} rekonstruiert die Architektur und die Evidence-Chain von \agora{}. Kapitel~\ref{ch:methodik} beschreibt die Audit-Pipeline. Kapitel~\ref{ch:ergebnisse} trägt die Befunde aus Code-Audit, Referenzlauf und Bewertung zusammen. Kapitel~\ref{ch:diskussion} diskutiert fünf Kernkontroversen. Kapitel~\ref{ch:schluss} beantwortet die Forschungsfrage und priorisiert den Ausblick. Anhang~\ref{ch:anhang-a} dokumentiert die Verifikations-Spot-Checks und die Methodik der Citation-Registry.

\section{Begriffliche Abgrenzung}\label{sec:begriffe}

Drei Begriffe werden im Folgenden streng in der hier definierten Bedeutung verwendet:

\begin{description}
  \item[Seed-Dokument] Ein vom Nutzer hochgeladenes Korpus-Dokument (PDF, Markdown oder Text), das als einzige \emph{echte externe Quelle} in die Pipeline eingeht. Alle anderen \code{EvidenceSourceKind}-Werte~\cite{adr-0013} sind entweder LLM-generiert, synthetisch oder abgeleitet.
  \item[Validierter Claim] Ein Claim, der eine \code{supports\_claim = True}-Zuordnung von der zweistufigen Entailment-Stufe erhalten hat. Diese Stufe kombiniert eine Cosine-Retrieval-Bewertung mit regelbasierten deterministischen Prüfungen~\cite{adr-0011}.
  \item[Provenance-Verlust] Ein Punkt in der Pipeline, an dem die Identität einer Information -- ihre doc-ID, chunk-ID oder Herkunftsquelle -- entweder wegfällt oder nicht in einen späteren Verarbeitungsschritt übernommen wird.
\end{description}

\section{Hinweis zum Reproduktionsstand}\label{sec:reproduktion}

Der Audit wurde am Repo-Stand \code{7e42ae34} durchgeführt. Zwischen Stand des Audits und dieser Niederschrift liegen zwei Merges in den Hauptzweig, die zentrale Befunde teils entschärfen: ADR-0013 Slice~1 Teil~A (PR~\#1155) und der Interview-Binding-Fix (Commit \code{d7d9f0a4}, PR~\#1151). Die Stellen sind in der Diskussion ausdrücklich markiert.

% =====================================================================
\part{Stand der Forschung}
% =====================================================================

% =====================================================================
% Kapitel 2: Stand der Forschung
% =====================================================================

\chapter{Stand der Forschung}\label{ch:related}

Der Stand der Forschung wird in sechs Themencluster gegliedert, die für die Bewertung von \agora{} direkt relevant sind. Pro Cluster werden zwei bis vier Schlüsselarbeiten referenziert und auf die Forschungslücke zugespitzt, die diese Arbeit schließen hilft.

\section{Synthetische Personas und generative Agenten}\label{sec:personas}

Die zentralen Vorarbeiten stammen von~\cite{park2023generative}, die in der Smallville-Simulation 25 generative Agenten über zwei Tage agieren lassen und qualitative Plausibilität, nicht aber populationsstatistische Repräsentativität nachweisen. \cite{argyle2023silicon} zeigen, dass LLM-Persona-Conditioning Antwortverteilungen replizieren kann, betonen aber die Prompt-Abhängigkeit. \cite{horton2023homosilicus} demonstriert qualitative Replikation ökonomischer Verhaltensweisen. \cite{uas2025} untersuchen systematisch die Replikation realer Survey-Daten und identifizieren persistenten Bias.

\textbf{Forschungslücke:} Es existiert kein kanonisches Protokoll, um LLM-Personas systematisch gegen reale Populationsstatistiken zu validieren~\cite{task-f-literature-note}. Übertragbarkeit auf DACH-Populationen ist eine offene empirische Frage.

\section{Multi-Agent-Debate und Self-Consistency}\label{sec:debate}

\cite{du2023multiagent} zeigen, dass Multi-Agent-Debate die Factuality gegenüber einzelnen Modellen verbessert, dass der Effekt jedoch bei großen Modellen sättigt. \cite{wang2022selfconsistency} liefern das methodische Fundament für Self-Consistency-Ansätze; \cite{consistency2024} relativieren deren Wert durch den Befund, dass Konsistenz nur bedingt mit Korrektheit korreliert. \cite{hiddenbench2025} zeigen, dass Kollektivreasoning am Hidden-Profile-Problem systematisch versagt (30,1\,\% vs.\ 80,7\,\% für einen Einzelagent mit vollständiger Information). \cite{mast2025} klassifizieren Fehlermodi in Multi-Agent-Systemen und stellen fest, dass Benchmark-Gewinne oft minimal sind.

\textbf{Forschungslücke:} Studien, die Multi-Agent-Debate \emph{mit aktiver Quellenverifikation} kombinieren, sind selten;~\cite{zenodo-agora} ist eine Ausnahme.

\section{GraphRAG und Retrieval-Grounding}\label{sec:graphrag}

\cite{edge2024graphrag} zeigen, dass GraphRAG bei globalen, query-fokussierten Zusammenfassungen Standard-RAG übertrifft. Inline-Zitationen sind, wie~\cite{citednotverified} nachweisen, häufig fehlerhaft oder irrelevant -- ein direkter Bezugspunkt zu \agoras{} Evidence-Gating~\cite{adr-0002}.

\textbf{Forschungslücke:} Vergleichende Evaluationen zwischen GraphRAG und Standard-RAG bei festem Budget sind rar; Effizienz-Varianten wie E\textsuperscript{2}GraphRAG oder Practical GraphRAG sind erst seit 2025/2026 verfügbar.

\section{Halluzinationserkennung und Kalibrierung}\label{sec:halluzination}

\cite{bail2024genai} formuliert die Warnung vor Homogenität, Halluzination und fehlender Validierung in der sozialwissenschaftlichen Anwendung generativer KI. \cite{consistency2024} zeigt, dass Selbstvertrauen und Korrektheit entkoppelt sind.

\textbf{Forschungslücke:} Verbindung von Halluzinationserkennung mit prozessualen Gating-Mechanismen -- wie \agora{} sie implementiert~\cite{adr-0002} -- ist methodisch unterentwickelt.

\section{Synthetische Daten und Model-Collapse}\label{sec:modelcollapse}

\cite{shumailov2024curse} zeigen, dass rekursives Training auf synthetisch generierten Daten die Verteilung verengt und schließlich zum Vergessen führt. Diese Erkenntnis ist für \agora{} zentral, da die Persona-Population vollständig synthetisch ist und über mehrere Simulations-Runden weiter aggregiert wird.

\textbf{Forschungslücke:} Langfristige Stabilität synthetischer Multi-Agent-Simulationen über Wochen oder Monate wird kaum systematisch gemessen.

\section{Analytische Flexibilität und Reproduzierbarkeit}\label{sec:flexibility}

\cite{analyticflexibility} untersuchen 252 Konfigurationen eines LLM-basierten Sozialmodells und zeigen, dass die Ergebnisse stark von Prompt, Sampling und Modellwahl abhängen -- ein P-Hacking-Äquivalent.

\textbf{Forschungslücke:} Standardisierte Reproduzierbarkeitsprotokolle für LLM-basierte Simulationen fehlen.

\section{Synthese: Wo setzt \agora{} an?}\label{sec:synthese}

\agora{} sitzt methodisch am Schnittpunkt von drei Clustern: (1)~\emph{Provenance-Grounding} durch serverseitige Anchor-Erzwingung, (2)~\emph{Anti-Self-Confirmation} durch Echo-Chamber-Index und Confidence-Capping, und (3)~\emph{Halluzinationsprävention} durch zweistufiges Binding. Damit berührt es Forschungslücken, die in der Literatur benannt sind, ohne sie bislang geschlossen zu haben. Der vorliegende Beitrag macht diese Lücken am konkreten Code sichtbar und formuliert eine Roadmap, die sie schließen könnte.

% =====================================================================
\part{System und Pipeline}
% =====================================================================

% =====================================================================
% Kapitel 3: System und Pipeline
% =====================================================================

\chapter{\agora{}: System und Pipeline}\label{ch:system}

\section{Architekturüberblick}\label{sec:architektur}

\agora{} ist eine lokal oder kontrolliert hybrid betreibbare Multi-Agent-Plattform für simulierte DACH-Zielgruppen-, Stakeholder- und Marktreaktionen. Die Plattform ist als Flask/Python-3.14-Backend mit Neo4j-Graph, Redis-Cache, OASIS- bzw.\ CAMEL-Simulations-Subprozess und Vue-3-Frontend strukturiert. Die Code-Basis umfasst zum Auditzeitpunkt etwa 3500~LOC im Backend-Service-Layer, 5500~LOC im OASIS-Subprozess und etwa 3900~LOC Graph-bezogenen Code. Der Reifegrad ist \emph{0.9.3 Stability Beta}; das Ziel ist nach \code{ROADMAP.md} eine stabile Single-User-Version \code{1.0.0}\footnote{\raggedright Status-Informationen aus~\code{AGENTS.md} und \code{ROADMAP.md} im Repository.}

Für die vorliegende Arbeit sind zwei Schichten zentral: die \textbf{Evidence-Pipeline} (vom Seed-Dokument zum Claim) und die \textbf{Simulations-Pipeline} (Persona-Population, OASIS-Aktionen, Echo-Index). Beide sind über die \code{EvidenceSourceKind}-Taxonomie~\cite{adr-0011} miteinander verbunden.

\section{Die Evidence-Chain}\label{sec:evidencechain}

Abbildung~\ref{fig:evidence-chain} rekonstruiert die Evidence-Chain vom Seed-Dokument bis zum ausgelieferten Claim in 17 Stufen. Die Rekonstruktion beruht auf statischer Code-Analyse von~\code{backend/app/services/} und~\code{backend/app/contracts/}, validiert durch die Code-Review-Graph-Werkzeuge. Sieben Stufen (\textsc{✗}) markieren code-verifizierte Provenance-Verlustpunkte.

\begin{figure}[ht]
  \centering
  % =====================================================================
% figures/evidence-chain.tex
% TikZ-Rekonstruktion der Evidence-Chain (Audit §1, vereinfacht)
% =====================================================================

\begin{tikzpicture}[
  font=\small,
  stage/.style={draw, rounded corners=2pt, minimum width=44mm, minimum height=9mm,
                align=center, fill=blue!5},
  losspoint/.style={draw, rounded corners=2pt, minimum width=44mm, minimum height=9mm,
                    align=center, fill=red!8, drop shadow},
  arrow/.style={-{Stealth[length=2.5mm]}, thick, draw=blue!60!black},
  lossyearrow/.style={-{Stealth[length=2.5mm]}, thick, draw=red!70!black, dashed},
  label/.style={font=\footnotesize, fill=white, inner sep=1pt},
  losslabel/.style={label, draw=red!70!black, rounded corners=1pt, inner sep=1.5pt}
]

% Hauptkette (Stufen 1-17)
\node[stage] (s1)  {[1] Seed-Dok.\\(Upload: PDF/MD/TXT)};
\node[stage, below=8mm of s1]  (s2)  {[2] Chunks\\\texttt{split\_text}};
\node[losspoint, below=8mm of s2]  (s3)  {\textbf{[3] KG}\\Neo4j \texttt{add\_text}};
\node[stage, below=8mm of s3]  (s4)  {[4] Retrieval\\\texttt{search\_graph}};
\node[losspoint, below=8mm of s4]  (s5)  {\textbf{[5] Report Tools}\\\texttt{\_record\_tool\_evidence}};
\node[stage, below=8mm of s5]  (s6)  {[6] Persona/Sim\\OASIS/CAMEL};
\node[stage, below=8mm of s6]  (s7)  {[7] Interview\\\texttt{AgentInterview}};
\node[stage, below=8mm of s7]  (s8)  {[8] Evid.-Norm.\\\texttt{\_TYPE\_TO\_SOURCE\_KIND}};
\node[stage, below=8mm of s8]  (s9)  {[9] Evid.-Ident.\\\texttt{build\_evidence\_id}};
\node[stage, below=8mm of s9]  (s10)  {[10] Evid.-Index\\\texttt{evidence-map.json}};
\node[stage, below=8mm of s10] (s11) {[11] Claim-Bind.\\\texttt{evidence\_binder}};
\node[stage, below=8mm of s11] (s12) {[12] Gating\\(ADR-0002)};
\node[stage, below=8mm of s12] (s13) {[13] Confidence\\\texttt{compute\_conf.}};
\node[stage, below=8mm of s13] (s14) {[14] Hypothesen\\\texttt{dedup\_and\_cap}};
\node[stage, below=8mm of s14] (s15) {[15] Red-Team\\(LLM-Judge)};
\node[losspoint, below=8mm of s15] (s16) {\textbf{[16] Renderer}\\\texttt{markdown\_renderer}};
\node[stage, below=8mm of s16] (s17) {[17] Export\\(JSON/ZIP/MD/CSV)};

% Pfeile
\draw[arrow] (s1) -- (s2);
\draw[lossyearrow] (s2) -- (s3) node[midway, losslabel, xshift=3mm] {\lossmark{} 1};
\draw[lossyearrow] (s3) -- (s4) node[midway, losslabel, xshift=3mm] {\lossmark{} 2,3,4};
\draw[lossyearrow] (s4) -- (s5) node[midway, losslabel, xshift=3mm] {\lossmark{} 5};
\draw[lossyearrow] (s5) -- (s6) node[midway, losslabel, xshift=3mm] {\lossmark{} 6,7};
\draw[arrow] (s6) -- (s7);
\draw[arrow] (s7) -- (s8);
\draw[arrow] (s8) -- (s9);
\draw[arrow] (s9) -- (s10);
\draw[arrow] (s10) -- (s11);
\draw[arrow] (s11) -- (s12);
\draw[arrow] (s12) -- (s13);
\draw[arrow] (s13) -- (s14);
\draw[arrow] (s14) -- (s15);
\draw[lossyearrow] (s15) -- (s16);
\draw[arrow] (s16) -- (s17);

% Legende (inline, ohne tikz-matrix)
\node[stage, minimum width=22mm, minimum height=5mm, below=10mm of s17, xshift=-30mm] (leg1) {};
\node[losspoint, minimum width=22mm, minimum height=5mm, right=2mm of leg1] (leg2) {};
\node[label, right=2mm of leg2] {\lossmark{} = Provenance-Verlustpunkt};

\end{tikzpicture}

  \caption{Rekonstruierte Evidence-Chain vom Seed-Dokument bis zum Claim. Die sieben mit \textsc{✗} markierten Stufen sind code-verifizierte Provenance-Verlustpunkte.}
  \label{fig:evidence-chain}
\end{figure}

\section{Die \code{EvidenceSourceKind}-Taxonomie}\label{sec:sourcekind}

Der Vertrag~\code{backend/app/contracts/report\_contract.py} definiert sechs Werte für den Enum~\code{EvidenceSourceKind}:

\begin{itemize}
  \item \code{seed\_corpus} -- echte externe Quelle aus dem hochgeladenen Korpus. Aktuell nur als Zielzustand in~\cite{adr-0013} spezifiziert; im Referenzlauf~\cite{audit-r2} nicht erreicht.
  \item \code{agent\_quote} -- Persona-Interview-Zitat. Provenance über~\code{persona\_id}, \code{stakeholder\_group}, \code{quote}.
  \item \code{agent\_action} -- sozialer Akt aus der Simulation (Post, Comment, Reshare). Provenance über \code{sim\_action\_log} (Plattform, Runde, Agent, Aktion, Zeitstempel).
  \item \code{graph\_relation} -- Fakt aus dem Neo4j-Graph. Aktuell nur syntaktisch referenzierbar; \emph{keine} Dokumentidentität.
  \item \code{web\_source} -- URL aus externer Recherche. URL im Anker; Fetch ist echtzeitabhängig und kann verfallen.
  \item \code{inferred} -- Standard-Fallback, wenn die Herkunft unbekannt ist. \code{Confidence}-Gewicht ist 0{,}0, d.\,h.\ ein Claim auf \code{inferred} endet als Hypothese.
\end{itemize}

Die Trennung ist die Grundlage für die nachfolgend beschriebene Pipeline -- und zugleich der Punkt, an dem die sechs Werte in der praktischen Wirksamkeit sehr ungleich verteilt sind.

\section{Zweistufiges Claim-Binding}\label{sec:binding}

Die \code{evidence\_binder} und \code{evidence\_entailment} (in \code{backend/app/services/}) implementieren ein zweistufiges Binding:

\begin{enumerate}
  \item \textbf{Retrieval-Stage:} Cosine-Ähnlichkeit zwischen Claim-Vektor und Evidence-Vektor. Schwellen 0{,}55 (niedrige Confidence) und 0{,}65 (hohe Confidence).
  \item \textbf{Entailment-Stage:} Regelbasierte deterministische Prüfung -- SUPPORTED, CONTRADICTED, RELATED\_ONLY oder INSUFFICIENT. Diese Stufe prüft Zahl, Bezugsgruppe, Menge und Modalität explizit vor dem Embedding und verhindert damit das in~\cite{adr-0011} belegte Defektmuster \enquote{61\,\% Zeitersparnis $\to$ 61\,\% bewerten positiv}.
\end{enumerate}

Das \code{supports\_claim}-Feld wird ausschließlich in der Entailment-Stufe gesetzt; dies ist im Code-Kommentar zu \code{evidence\_entailment.py} ausdrücklich festgehalten. Diese Architektur ist zentral für das Verständnis des gesamten Audit-Befundes, weshalb sie in Kapitel~\ref{ch:ergebnisse} erneut aufgegriffen wird.

\section{Anti-Self-Confirmation-Mechanismen}\label{sec:anticonf}

\agora{} implementiert vier strukturelle Anti-Self-Confirmation-Mechanismen, die in der folgenden Tabelle zusammengefasst sind. Ausführlich werden sie in Kapitel~\ref{ch:ergebnisse} bewertet.

\begin{table}[ht]
  \centering
  \small
  \begin{tabular}{@{}p{0.32\textwidth}p{0.58\textwidth}@{}}
    \toprule
    \textbf{Mechanismus} & \textbf{Wirkung} \\
    \midrule
    Echo-Chamber-Index (\code{network\_analytics.py}) &
    Misst intra-cluster-Anteil synthetischer Aktionen; drosselt Confidence bei Werten $> 0{,}75$ via \code{apply\_echo\_cap}. \\
    \midrule
    DACH-Quoten-Kalibrierung &
    Verhindert willkürliche Persona-Population; IT-Anteil hard-gecappt $\leq 12\,\%$; DACH-Branchenanteile gegen Destatis WZ 2008 / Mikrozensus 2024 kalibriert. \\
    \midrule
    Wording-Glossar (\code{\_red\_team\_system\_prompt}) &
    Verbietet Vorhersagesprache; Volltext-Snapshot; Red-Team-Enforcement post-Assembly. \\
    \midrule
    Confidence-Capping &
    Caps für $< 2$ Evidence-Items (0{,}59), \code{model\_generated\_inference}-Gewicht 0{,}0; Echo-Cap auf 0{,}84/medium. \\
    \bottomrule
  \end{tabular}
  \caption{Anti-Self-Confirmation-Mechanismen in \agora{}; Wirkung und Code-Anker.}
  \label{tab:anticonf}
\end{table}

\section{Simulations-Subprozess}\label{sec:simulation}

Der Simulations-Subprozess integriert OASIS- bzw.\ CAMEL-basierte soziale Simulation. Eine Population synthetischer Personas wird über bis zu 5500~LOC Subprozess-Code über mehrere Runden hinweg sozial aktiv (Posts, Comments, Reshares). Der Echo-Chamber-Index berechnet sich als Verhältnis der Intra-Cluster-Aktionen zur Gesamtzahl und koppelt an das Confidence-Gating. Die Simulation ist \emph{code-verifiziert nicht deterministisch} (\code{random.*} ohne \code{random.seed()} an mehreren Stellen; einzige Determinismus-Insel ist \code{louvain(seed=42)} im Community-Detection-Schritt).

\section{Vom Code zum Paper: Annotationsstrategie}\label{sec:annotations}

Jede Code-Aussage in den folgenden Kapiteln ist über einen der vier Schlüssel\texttt{\code{C<n>}}, \texttt{\code{A<n>}}, \texttt{\code{I<n>}} und \texttt{\code{R<n>}} belegt. Eine vollständige Tabelle der Belege ist in der Citation-Registry~\cite{citation-registry} hinterlegt; die Datei listet 57~Code-Belege, 5~ADRs, 12~Issues, 4~Referenzläufe und 14~Literaturschlüssel. Im Text werden die Code-Belege als \texttt{\src{C11}{split\_text ohne Manifest, \code{graph\_build.py:439}}} inline nachgewiesen; \code{AGENTS.md} und \code{README.md} werden ausdrücklich \emph{nicht} als Belege herangezogen.

% =====================================================================
\part{Methodik}
% =====================================================================

\chapter{Methodik}\label{ch:methodik}

\section{Übersicht: Audit-Pipeline}\label{sec:audit-pipeline}

Die Arbeit folgt einer strukturierten Audit-Pipeline, die statische Quellcode-Analyse mit normativer Bewertung verknüpft. Der Lead-Autor koordiniert acht parallele Subagenten (Tasks~A–G, ergänzt um B1–B3 zur Sim-Tiefe); jeder Subagent produziert eine Research-Note unter \code{docs/paper/research-notes/}. Der Lead führt Counter-Review, Citation-Registry und Schlussredaktion durch. Die Pipeline ist in der Begleitdokumentation des Repositories vollständig dokumentiert~\cite{zenodo-agora}.

\begin{figure}[ht]
  \centering
  \begin{tikzpicture}[>=Stealth, node distance=8mm and 6mm]
    \tikzstyle{box}=[draw, rounded corners, minimum width=32mm, minimum height=8mm, align=center, font=\small]
    \tikzstyle{arr}=[->, thick]
    \node[box] (lead) {Lead\\(Plan, Registry, Review)};
    \node[box, right=of lead] (a) {Task A\\Ingest \& Graph};
    \node[box, right=of a] (b) {Task B1--B3\\Simulation};
    \node[box, right=of b] (c) {Task C\\Report \& Gating};
    \node[box, below=of a] (d) {Task D\\Provenance};
    \node[box, below=of b] (e) {Task E\\Architektur};
    \node[box, below=of c] (f) {Task F\\Literatur};
    \node[box, below=of d] (g) {Task G\\Tests};
    \draw[arr] (lead) -- (a);
    \draw[arr] (lead) -- (b);
    \draw[arr] (lead) -- (c);
    \draw[arr] (a) -- (d);
    \draw[arr] (b) -- (e);
    \draw[arr] (c) -- (f);
    \draw[arr] (d) -- (g);
  \end{tikzpicture}
  \caption{Subagent-Topologie der Audit-Pipeline. Der Lead bleibt im Kontrollthread; jeder Subagent arbeitet isoliert und liefert eine Note zurück.}
  \label{fig:audit-pipeline}
\end{figure}

\section{Code-Review-Graph als Strukturspeicher}\label{sec:crg}

Das zentrale Werkzeug ist der im Repository gepflegte Code-Review-Graph (CRG). Er parst den Quellcode mit Tree-Sitter und bildet einen Graphen aus Dateien, Klassen, Funktionen und Aufrufbeziehungen. Die folgenden Operationen wurden eingesetzt:

\begin{itemize}
  \item \code{semantic\_search\_nodes\_tool} -- zur Identifikation relevanter Symbole zu einem Konzept (z.\,B.\ \enquote{evidence binding}, \enquote{claim entailment}).
  \item \code{query\_graph\_tool} -- mit dem Pattern~\code{callers\_of} bzw.~\code{callees\_of} zur Extraktion der Aufrufkette.
  \item \code{get\_impact\_radius\_tool} -- zur Eingrenzung des Initial-Scan auf ~5–10 Knoten pro Subagent.
  \item \code{detect\_changes\_tool} -- zur Verifikation der Befunde gegen den letzten Commit auf dem Branch.
\end{itemize}

Diese Vorgehensweise reduziert die Kontext-Last jedes Subagenten auf -- nach interner Messung -- etwa 30–40\,\% im Vergleich zu unkontrolliertem \code{grep}/\code{Read}. Der Lead behält die Kontext-Übersicht und führt die Synthese durch.

\section{Verifikation am Code}\label{sec:verifikation}

Jede Code-Aussage wurde am Branch-Stand \code{7e42ae34} verifiziert. Verwendet wurden:

\begin{itemize}
  \item Direktes Lesen der Quellen via \code{Read} (nur für die durch CRG benannten ~5–10 Stellen je Subagent).
  \item Verarbeitende Skripte via \code{ctx\_execute\_file} -- insbesondere für das Parsen der \code{evidence\_map.json}-Schemas, der Verifikations-Spot-Checks und des Failure-Mode-Patterns.
  \item Statische Aufrufanalyse via \code{query\_graph\_tool} für die Rekonstruktion der Evidence-Chain.
\end{itemize}

Ein zentrales methodisches Prinzip: \emph{kein} Verlassen auf~\code{README.md}, \code{docs/architecture.md} oder \code{VISION.md} als Belegquelle. Diese Dateien werden als Sekundärbeschreibungen behandelt; Aussagen aus ihnen fließen nur ein, wenn sie am Code bestätigt werden.

\section{Counter-Review}\label{sec:counter-review}

Vor der Schlussfassung wurde die zentrale These (s.\ Kapitel~\ref{ch:schluss}) einem Counter-Review nach~\cite{arxiv:counter-review-methodik} unterzogen. Die folgenden Einwände wurden geprüft:

\begin{enumerate}
  \item \textbf{Könnte die Schlussfolgerung falsch sein?} Die These \enquote{0~validierte Claims = kein Mehrwert} würde falsch, wenn R2 ein Einzelfall wäre. Gegenbeleg: R1~\cite{audit-r1} (historisch) zeigt 25~Claims auf 100\,\% inferred, ADR-0011~\cite{adr-0011} bestätigt strukturell. Eingeschränkt nach Sim-Verifikation: R2 dokumentiert den Pre-Fix-Stand; ein E2E-Re-Run gegen \code{d7d9f0a4} könnte die Claim-Zahl erhöhen -- die These ruht damit auf einer Säule weniger, ist aber nicht umgestoßen.
  \item \textbf{Single-Source-Abhängigkeit?} These~D (OASIS-Mehrwert ungeprüft) ruht allein auf Task~E; abgestützt durch L10, L11.
  \item \textbf{Stale Quellen?} Alle Code-Belege \code{AS\_OF}~2026-08-09. ADR-0011 (13~Tage alt) nach manueller Prüfung weiterhin gültig.
  \item \textbf{$\geq 3$ Probleme gefunden:} 5~Counter-Claims geprüft, alle gehalten oder als Negativbeweis markiert.
\end{enumerate}

\section{Reproduzierbarkeit der Methode}\label{sec:reproduzierbarkeit}

Die Methode ist in drei Hinsichten reproduzierbar:

\begin{itemize}
  \item \textbf{Pipeline:} Die Reihenfolge Lead $\to$ Subagenten $\to$ Notes $\to$ Registry $\to$ Counter-Review ist im Repository dokumentiert (s.\ Task-F, \code{notes/methodik.md}).
  \item \textbf{Belege:} Alle 33~Belege sind in der Citation-Registry~\cite{citation-registry} mit Symbol, Datei, Zeilennummer verankert. Zeilennummern können bei Refactors driften; Symbole bleiben stabil.
  \item \textbf{Software-Artefakt:} Ein eingefrorener Snapshot des Repository-Standes ist unter der DOI \code{10.5281/zenodo.21855023} hinterlegt~\cite{zenodo-agora}.
\end{itemize}

\section{Einschränkungen}\label{sec:einschraenkungen}

Drei Einschränkungen sind zu nennen:

\begin{enumerate}
  \item \textbf{Pre-Fix-Stand:} Der zentrale Referenzlauf~\cite{audit-r2} wurde vor dem Merge von~\code{d7d9f0a4} (Interview-Binding-Fix) erzeugt. Die Befunde sind dadurch konservativ; ein Re-Run könnte die Zahl validierter Claims erhöhen.
  \item \textbf{Negativbeweis für OASIS-Mehrwert:} Die Aussage \enquote{OASIS-Mehrwert ungeprüft} stützt sich auf eine leere Suche nach \code{ground\_truth}, \code{destatis}, \code{census} oder \code{sim\_validity} in \code{backend/app/}. Das ist ein Negativbeweis -- kein Beweis, dass der Mehrwert nicht existiert, sondern dass er nicht operationalisiert ist.
  \item \textbf{Code-Audit $\neq$ Empirie:} Die Befunde sind statische Quellcode-Analysen und Output-Inspektionen. Aussagen über die \emph{Wirkung} der Mechanismen auf reale Populationen können daraus nicht abgeleitet werden; das wäre Aufgabe einer prospektiven Sim-Validity-Studie.
\end{enumerate}

\section{Annotation der Belege}\label{sec:annotations2}

Im Text werden Code-Belege als \texttt{\src{C11}{split\_text ohne Manifest, \code{graph\_build.py:439}}} inline ausgewiesen. Die Notation folgt~\cite{citation-registry}; eine ausführliche Belegliste ist im Anhang dokumentiert.

% =====================================================================
\part{Ergebnisse}
% =====================================================================

% =====================================================================
% Kapitel 5: Ergebnisse
% =====================================================================

\chapter{Ergebnisse}\label{ch:ergebnisse}

\section{Referenzlauf und Baseline}\label{sec:r2}

Der Referenzlauf~\cite{audit-r2} (2026-08-09, Eingabe: 16 Seed-Dokumente, 5~Stakeholder-Gruppen, 12~Personas) erzeugte \textbf{null validierte Claims}. Der gesamte Inhalt wanderte in das Hypothesen-Slot~\src{C36}{dedup\_and\_cap\_hypotheses}. Der historische Vergleichslauf~\cite{audit-r1} (2026-07-27) erzeugte 25~Claims, alle~\code{medium}, auf 125~Evidence-Items, die zu 100\,\% den Wert~\code{inferred} trugen -- der dokumentierte Vor-Fix-Zustand aus~\cite{adr-0011}.

\textbf{Beobachtung:} Der Sprung von 25~medium-Claims (mit kaputtem Binding) zu 0~validierten Claims (mit korrektem Binding) belegt, dass das Gating funktioniert -- aber ins Leere greift, weil der Upstream keine kanonische Evidence liefert.

\section{10-Dimensionen-Bewertung}\label{sec:scoring}

Die folgende Bewertung folgt dem Schema aus dem zugrundeliegenden Audit~\cite{zenodo-agora} und ist die Grundlage der Gesamtscores in Kapitel~\ref{ch:diskussion}. Jede Dimension wird auf einer Skala von 0 (v{\"o}llig unzureichend) bis 10 (vorbildlich) bewertet.

\begin{table}[ht]
  \centering
  \small
  \begin{tabular}{@{}p{0.04\textwidth}p{0.32\textwidth}p{0.05\textwidth}p{0.45\textwidth}@{}}
    \toprule
    \textbf{\#} & \textbf{Dimension} & \textbf{Score} & \textbf{Begründung} \\
    \midrule
    1 & Provenance-Architektur (Konzept) & 8 & Sechs source\_kind, Default \code{inferred}, Anker-Konzept, Identit{\"a}tswechsel; konzeptionell stark und durch ADRs verankert. \\
    2 & Provenance-Umsetzung (Ist) & 3 & Teil~B fehlt, Anker opak, \code{medium} unerreichbar, 0~validierte Claims. \\
    3 & Evidence-Binding/Rigour & 7 & Zweistufig, Determinismus vor Embedding, Modalit{\"a}tstrennung; greift ins Leere ohne Upstream-Evidence. \\
    4 & Confidence-Honesty & 7 & Mehrkomponenten, Caps, contradiction-penalty; aber Konzeptvermischung (key\_takeaways high bei claims:[]). \\
    5 & Anti-Self-Confirmation & 7 & Echo-Cap, Red-Team, Wording-Gate; aber Schwellen nicht kalibriert, nur cross-stakeholder. \\
    6 & Testabdeckung & 5 & 363~Tests, starke Contract-/Gating-Tests; aber kein E2E-Provenance, kein Baseline, kein Reproducibility-Test. \\
    7 & Reproduzierbarkeit & 2 & Simulation nicht deterministisch; E2E nur Stub; kein Same-Seed-Same-Report-Test. \\
    8 & Sim-Validity/Empirie & 2 & Keine Metrik, kein Ground-Truth, kein A/B; Referenzlauf inhaltlich nicht vertrauensw{\"u}rdig. \\
    9 & Komplexit{\"a}ts{\"o}konomie & 5 & Neo4j/Redis/CAMEL f{\"u}r Single-User overkill; OASIS \textasciitilde5500~LOC ohne Validierung; Evidence-Pipeline hingegen gerechtfertigt. \\
    10 & Doku/ADR-Honesty & 9 & ADR-0011 dokumentiert eigenen Fehler offen; 5~Hartanker prozessual gesch{\"u}tzt; Roadmap nennt Schw{\"a}chen beim Namen. \\
    \bottomrule
  \end{tabular}
  \caption{10-Dimensionen-Bewertung von \agora{} zum Auditzeitpunkt.}
  \label{tab:scoring}
\end{table}

\section{Failure-Mode-Analyse}\label{sec:failures}

Aus dem Audit wurden 15~Failure-Modes (FM) extrahiert, geordnet nach Severity. Tabelle~\ref{tab:failures} zeigt die kritischen Befunde; die vollständige Liste findet sich im Anhang.

\begin{table}[ht]
  \centering
  \small
  \begin{tabular}{@{}p{0.04\textwidth}p{0.32\textwidth}p{0.10\textwidth}p{0.44\textwidth}@{}}
    \toprule
    \textbf{FM} & \textbf{Failure-Mode} & \textbf{Severity} & \textbf{Code-Beleg / Gegenmaßnahme} \\
    \midrule
    F1 & Seed-Provenance end-to-end gebrochen (Teil~B fehlt) & Critical & \src{C11}{split\_text ohne Manifest}; ADR-0013 Teil~B \\
    F2 & Interview$\to$Evidence-Binding-Defekt; Fix in HEAD (\code{d7d9f0a4}), E2E-Re-Run offen & Medium & \src{C50}{producer\_key + quote + persona\_stakeholder\_group} \\
    F3 & Opaker \code{seed\_doc:}-Anker (LLM nennt eigene Quelle, niemand pr{\"u}ft nach) & Critical & \src{C9}{Anker opak}, \src{C10}{kein Lookup} \\
    F4 & Gate zu streng f{\"u}r eigenen Throughput $\to$ 0~validierte Claims & High & Folge von F1+F2 \\
    F5 & Graph-Fakten ohne Dokumentidentit{\"a}t (\code{List[str]}) & High & \src{C19}{Episode ohne doc\_id}, \src{C21}{kein source\_id\_anchor} \\
    F6 & \code{medium}-Stufe aktuell unerreichbar & High & \src{C8}{has\_agent\_grounded\_evidence} \\
    F14 & Neo4j f{\"u}r Single-User-Local overkill (\textasciitilde3900~LOC Graph-Code) & Low & \src{C40}{Komplexit{\"a}t} \\
    F16 & \code{simulated\_hours} nicht in API/Frontend exponiert & Low & Issue~1018~\cite{issue-1018} \\
    \bottomrule
  \end{tabular}
  \caption{Auswahl kritischer Failure-Modes (vollständige Liste im Anhang).}
  \label{tab:failures}
\end{table}

\section{Provenance-Verlustpunkte}\label{sec:loss}

Die Rekonstruktion der Evidence-Chain (Abbildung~\ref{fig:evidence-chain}) identifiziert sieben code-verifizierte Provenance-Verlustpunkte. Sie sind in Tabelle~\ref{tab:loss} zusammengefasst.

\begin{table}[ht]
  \centering
  \small
  \begin{tabular}{@{}p{0.05\textwidth}p{0.45\textwidth}p{0.40\textwidth}@{}}
    \toprule
    \textbf{\#} & \textbf{Verlustpunkt} & \textbf{Beleg} \\
    \midrule
    \lossmarkX{1} & \code{split\_text} ohne Manifest; \code{List[str]} ohne \code{doc\_id/chunk\_id} & \src{C11}{graph\_build.py:439} \\
    \lossmarkX{2} & \code{episode\_id = uuid4()} ohne Dateibezug & \src{C15}{} \\
    \lossmarkX{3} & \code{Episode.data} = roher Chunk-String, kein \code{document\_id} & \src{C16}{} \\
    \lossmarkX{4} & \code{RELATION.episode\_ids = [uuid]}, keine Dokument-Referenz & \src{C17}{} \\
    \lossmarkX{5} & \code{SearchResult.facts = List[str]}, keine Herkunft & \src{C19}{}, \src{C20}{} \\
    \lossmarkX{6} & Graph-Fakt $\to$ \code{source\_kind=graph\_relation}, kein \code{source\_id\_anchor} & \src{C21}{}, \src{C7}{} \\
    \lossmarkX{7} & \code{seed\_doc:}-Pr{\"a}fix opak akzeptiert, kein Lookup & \src{C9}{}, \src{C10}{} \\
    \bottomrule
  \end{tabular}
  \caption{Sieben code-verifizierte Provenance-Verlustpunkte in der Evidence-Chain.}
  \label{tab:loss}
\end{table}

\section{Prompt-vs-\agora{}-Matrix}\label{sec:promptvsagora}

Die Matrix prüft für jeden \agora{}-Mechanismus, ob er durch einen einzelnen starken LLM-Prompt ersetzbar w{\"a}re. Bewertung: \lossmark=nein (strukturell), $\bullet$=teilweise (per Instruktion erreichbar), \okmark=ja.

\begin{table}[ht]
  \centering
  \small
  \begin{tabular}{@{}p{0.40\textwidth}p{0.05\textwidth}p{0.45\textwidth}@{}}
    \toprule
    \textbf{Mechanismus} & \textbf{Prompt?} & \textbf{Bemerkung} \\
    \midrule
    Echo-Chamber-Index + \code{apply\_echo\_cap} & \lossmark & Strukturelle Messung; nicht prompt-basierbar \\
    Wording-Glossar / Anti-Prediction & $\bullet$ & Per Prompt approximierbar, aber Snapshot-gepinnt \\
    Zweistufiges Binding (Retrieval vs.\ Entailment) & \lossmark & Zwei verschiedene Fragen, nicht ein Prompt \\
    Deterministische Checks vor Embedding & \lossmark & Strukturell, nicht prompt-basierbar \\
    \code{source\_kind}-Trennung (6~Stufen) & \lossmark & Trennt im Datentyp, nicht im Prompt \\
    Kanonische \code{evidence\_id} (SHA-256) & \lossmark & Reproduzierbare Identit{\"a}t \\
    Red-Team-Review (LLM-Judge post-Assembly) & $\bullet$ & Self-Critique per Prompt m{\"o}glich, aber an \code{echo\_index} gekoppelt \\
    Default \code{inferred} (unbekannt = abgeleitet) & \lossmark & Im Vertrag erzwungen \\
    Mode-Routing (strict dropt, balanced {\"u}berspringt) & $\bullet$ & Im Contract-Validator \\
    Graceful Degradation (gate-decision-log vs.\ degradation-log) & \lossmark & Getrenntes Audit-Log \\
    Multi-Agent-Simulation (OASIS/CAMEL) & \lossmark & Mehrwert ungepr{\"u}ft, keine Sim-Validity-Metrik \\
    Downgrade als Identit{\"a}tswechsel & \lossmark & Im Hash erzwungen \\
    ADR-0002 Hartanker (5, unantastbar) & \lossmark & Prozessuale Garantie \\
    \bottomrule
  \end{tabular}
  \caption{Prompt-vs-\agora{}-Matrix: Welche Mechanismen sind strukturell (\lossmark) und damit nicht durch einen st{\"a}rkeren System-Prompt ersetzbar?}
  \label{tab:promptvsagora}
\end{table}

\section{Novelty-Matrix}\label{sec:novelty}

Zur Einordnung der konzeptionellen Neuheit wird eine vier-stufige Skala verwendet:

\begin{description}
  \item[L0] Standard-LLM-Prompt
  \item[L1] strukturierte Output-Felder
  \item[L2] typisierte Contracts/Validatoren
  \item[L3] mehrstufige deterministische Checks vor/statt Embedding
  \item[L4] konzeptionell neu in der Kombination
\end{description}

Tabelle~\ref{tab:novelty} ordnet die zentralen \agora{}-Eigenschaften ein.

\begin{table}[ht]
  \centering
  \small
  \begin{tabular}{@{}p{0.50\textwidth}p{0.05\textwidth}p{0.35\textwidth}@{}}
    \toprule
    \textbf{Eigenschaft} & \textbf{Level} & \textbf{Beleg} \\
    \midrule
    Report als Textausgabe & L0 & Standard \\
    Strukturierte Claim/Evidence-Felder & L1 & Standard f{\"u}r Agent-Frameworks \\
    Pydantic-Contracts, \code{EvidenceSourceKind}-Enum, Gating-Validatoren & L2 & Standard in typisierten Backends \\
    Zweistufiges Binding & L3 & ADR-0011 motiviert durch 7~belegte Defekte \\
    Deterministische Checks vor Embedding & L3 & Verhindert Zahlenwanderung \\
    Server-seitiger Dokument-Anker (Konzept) & L3 / L0 & Konzept: L3, Umsetzung: L0 \\
    Echo-Chamber-Index + \code{apply\_echo\_cap} & L4 & Kopplung synthetischer Homogenit{\"a}tsmessung an Confidence-Cap \\
    Wording-Glossar als Snapshot + Red-Team-Enforcement & L3 & Prozessual verankert \\
    Downgrade als Identit{\"a}tswechsel & L4 & Erzwingt atomare Umschl{\"u}sselung \\
    \bottomrule
  \end{tabular}
  \caption{Novelty-Matrix: Einordnung zentraler \agora{}-Eigenschaften auf den Stufen~L0--L4.}
  \label{tab:novelty}
\end{table}

\section{Anti-Self-Confirmation-Bewertung}\label{sec:anticonfeval}

Aus der Forschungsbasis~\cite{consistency2024,citednotverified,shumailov2024curse,hiddenbench2025} ergibt sich die Frage, welche Anti-Self-Confirmation-Mechanismen \agora{} tats{\"a}chlich implementiert. Tabelle~\ref{tab:anticonfeval} bewertet die Mechanismen gegen{\"u}ber den in der Literatur dokumentierten Risiken.

\begin{table}[ht]
  \centering
  \small
  \begin{tabular}{@{}p{0.36\textwidth}p{0.34\textwidth}p{0.20\textwidth}@{}}
    \toprule
    \textbf{Mechanismus} & \textbf{Wirkung gegen das Literatur-Risiko} & \textbf{Caveat} \\
    \midrule
    Echo-Chamber-Index & Misst synthetische Homogenit{\"a}t & Misst nicht reale Populationsvalidit{\"a}t \\
    DACH-Quoten-Kalibrierung & Verhindert willk{\"u}rkliche Persona-Population & Konstanten, kein Mikrodaten-Fit; LLM-Echttest CI-excluded \\
    \code{apply\_echo\_cap} & Drosselt polarisierte Echo-Kammern & Greift nur bei cross-stakeholder; Schwellen hartkodiert \\
    Red-Team-Review (LLM-Judge) & Pr{\"u}ft Claims auf Schw{\"a}chen & LLM-Judge kann Confirmation-Bias haben; \code{BudgetExceededError} durchgereicht~\cite{issue-978} \\
    Wording-Glossar & Positioniert als Szenarienanalyse & An \code{echo\_index} gekoppelt, nicht an Claim-Falschheit \\
    Default \code{inferred} & Verhindert \enquote{jedes Item = Dokumentfakt} & \code{inferred}-Gewicht 0{,}0 $\to$ Claim endet als Hypothese \\
    Zweistufiges Binding & Verhindert \enquote{Cosine-only supports\_claim=True} & Schwaches Embedding $\to$ Claim wird Hypothese \\
    \bottomrule
  \end{tabular}
  \caption{Bewertung der Anti-Self-Confirmation-Mechanismen gegen{\"u}ber den in der Literatur dokumentierten Risiken.}
  \label{tab:anticonfeval}
\end{table}

\textbf{Befund:} \agora{} hat \emph{konkrete, code-level Mechanismen}, die ein Single-Prompt nicht besitzt. Sie sind der st{\"a}rkste Teil der Architektur -- greifen jedoch ins Leere, weil der Upstream keine kanonische Evidence liefert (R2).

\section{Referenzl{\"a}ufe im Vergleich}\label{sec:referenzlaeufe}

\begin{table}[ht]
  \centering
  \small
  \begin{tabular}{@{}p{0.10\textwidth}p{0.18\textwidth}p{0.18\textwidth}p{0.18\textwidth}p{0.18\textwidth}@{}}
    \toprule
    \textbf{Lauf} & \textbf{Datum} & \textbf{Claims} & \textbf{Evidence-Items} & \textbf{Anteil \code{inferred}} \\
    \midrule
    R1~\cite{audit-r1} & 2026-07-27 & 25 (\emph{alle medium}) & 125 & 100\,\% \\
    R2~\cite{audit-r2} & 2026-08-09 & 0 (validiert) & 65 & 78\,\% \\
    R4~\cite{audit-r4} & 2026-07-25 & (report\_d9023bd1f55a) & -- & -- \\
    \bottomrule
  \end{tabular}
  \caption{Vergleich der Referenzl{\"a}ufe. R1 dokumentiert den Vor-Fix-Zustand; R2 den Pre-Fix-Stand des Interview-Binding-Defekts; R4 ist der in ADR-0011 als Beleg zitierte Report.}
  \label{tab:referenzlaeufe}
\end{table}

\section{Zwischenfazit}\label{sec:zwischenfazit}

Die Ergebnisse erlauben eine erste Antwort auf die Forschungsfrage: \textbf{Aktuell liefert \agora{} keinen messbaren Erkenntnisgewinn gegen{\"u}ber einem Single-Prompt}. Die Strukturen sind ehrlich konstruiert; sie greifen jedoch ins Leere, weil die Seed-Provenance end-to-end gebrochen ist (F1) und die Anker opak akzeptiert werden (F3). Diese Aussage wird in Kapitel~\ref{ch:diskussion} anhand von f{\"u}nf Kernkontroversen eingeordnet.

% =====================================================================
\part{Diskussion}
% =====================================================================

% =====================================================================
% Kapitel 6: Diskussion
% =====================================================================

\chapter{Diskussion}\label{ch:diskussion}

\section{Fünf Kernkontroversen}\label{sec:kontroversen}

Die Ergebnisse werfen fünf Kernkontroversen auf, die im Folgenden einzeln diskutiert werden.

\subsection{Kontroverse~1: Evidence-Pipeline als Differenzierer?}\label{sec:k1}

Die Pipeline wird in~\cite{adr-0011} als zentraler Mehrwert gegen{\"u}ber Standard-Agent-Frameworks ausgewiesen. Der Code belegt die strukturelle Eigenst{\"a}ndigkeit~\src{C24}{zweistufiges Binding}, \src{C25}{Determinismus vor Embedding}. Der Referenzlauf~\cite{audit-r2} zeigt jedoch null validierte Claims -- ein Hinweis darauf, dass die Pipeline aktuell nicht greift. \textbf{These:} konzeptionell echter Differenzierer, aktuell wirkungslos.

\subsection{Kontroverse~2: Anti-Self-Confirmation {\"u}ber Prompt-Level?}\label{sec:k2}

Echo-Cap~\src{C26}{} und Wording-Gate~\src{C30}{} sind strukturell, nicht prompt-basierbar. Aber: Schwellen hartkodiert (0{,}75 bzw.\ 0{,}84), nur cross-stakeholder, ohne externe Ground-Truth bleibt der Mechanismus ein Selbsttest~\cite{analyticflexibility,uas2025}. Die Mechanismen sind real, ihre \emph{Wirkung} ist nicht kalibriert.

\subsection{Kontroverse~3: 0~validierte Claims = kein Mehrwert?}\label{sec:k3}

Das Gate ist korrekt -- der Defekt liegt im Upstream (F1--F3).~\cite{adr-0013} dokumentiert explizit, dass unmittelbar nach Umsetzung von Teil~B \enquote{Reports schlechter aussehen werden als heute}. Der kurzfristige Qualit{\"a}tsverlust ist der Preis f{\"u}r {\"u}berpr{\"u}fbare Provenance. \textbf{These:} kein aktueller Mehrwert, aber das Gate ist nicht schuld.

\subsection{Kontroverse~4: OASIS-Mehrwert?}\label{sec:k4}

Oasis umfasst etwa 5500~LOC Subprozess-Code, ohne dass eine Sim-Validity-Metrik existiert~\cite{mast2025,hiddenbench2025}. Der Mehrwert der Simulation ist weder widerlegt noch best{\"a}tigt -- die zentrale unbewiesene Annahme des Systems.

\subsection{Kontroverse~5: ADR-Honesty = Ersatz f{\"u}r Mehrwert?}\label{sec:k5}

Die ehrliche Selbstkritik in ADR-0011~\cite{adr-0011} ist bemerkenswert und als Engineering-Leistung ernst zu nehmen. Sie ersetzt jedoch nicht den empirischen Nachweis eines Mehrwerts. Dokumentierte Fehler sind ein notwendiger, nicht hinreichender Schritt zur Vertrauensw{\"u}rdigkeit.

\section{Engineering-Score vs.\ Evidence-Research-Score}\label{sec:scores}

Aus der 10-Dimensionen-Bewertung (Tabelle~\ref{tab:scoring}) leiten sich zwei Gesamtscores ab, die das zentrale Spannungsfeld des Systems benennen:

\begin{align*}
  \mathrm{Engineering\text{-}Score} &= \frac{8 + 7 + 7 + 7 + 9}{5} = 7.6 \to 7.0 \\
  \mathrm{Evidence\text{-}Research\text{-}Score} &= \frac{3 + 5 + 2 + 2 + 5}{5} = 3.4 \to 3.5
\end{align*}

Die Scores messen Verschiedenes: der Engineering-Score bewertet die Konstruktion (Konzept, Contracts, Gating, Mechanismen, Dokumentation), der Evidence-Research-Score bewertet den nachweisbaren empirischen Mehrwert. \agora{} ist solide konstruiert (7{,}0), aber empirisch unbewiesen (3{,}5).

\textbf{Interpretation:} Die Diskrepanz ist \emph{nicht} als Schw{\"a}che, sondern als \emph{Arbeitsauftrag} zu lesen. Das System ist eine sorgf{\"a}ltig gebaute Maschine, die beweist, dass sie nichts beweisen kann -- und das ist, paradoxerweise, genau die Ehrlichkeit, die der Forschung zu synthetischen Personas fehlt~\cite{citednotverified,analyticflexibility}.

\section{Warum greift das Gate aktuell ins Leere?}\label{sec:gate-leer}

Das Evidence-Gating~\cite{adr-0002} ist deterministisch und auditierbar (\code{gate\_decision\_log}, \code{degradation\_log}). Es dropt oder degradiert Claims zuverl{\"a}ssig, wenn die Voraussetzungen nicht erf{\"u}llt sind. Im Referenzlauf~\cite{audit-r2} f{\"u}hrt das zu null validierten Claims -- nicht weil das Gate zu streng ist, sondern weil der Upstream (Seed-Ingest, Graph-Aufbau, Tool-Evidence-Erfassung) keine kanonische Evidence liefert. Die Konsequenz: das System erstickt an seiner eigenen Strenge.

Dieser Befund ist f{\"u}r jedes andere System mit {\"a}hnlicher Architektur relevant: \emph{ohne} nachgewiesene Provenance-Upstream f{\"u}hrt striktes Downstream-Gating zu null validierter Ausgabe.

\section{Limitationen}\label{sec:limitationen}

\begin{itemize}
  \item \textbf{Code-Audit $\neq$ Empirie.} Die Befunde sind statische Quellcode-Analyse und Output-Inspektion. {\"U}ber die Wirkung der Mechanismen auf reale Populationen kann diese Arbeit keine Aussage treffen.
  \item \textbf{Pre-Fix-Stand:} R2 wurde vor~\code{d7d9f0a4} erzeugt; ein Re-Run k{\"o}nnte die Claim-Zahl erh{\"o}hen -- die Konzepttrennung (F8) bleibt davon unber{\"u}hrt.
  \item \textbf{Negativbeweis f{\"u}r OASIS-Mehrwert:} Eine leere Suche nach \code{ground\_truth}-Anchern belegt \emph{nicht}, dass der Mehrwert nicht existiert, sondern dass er nicht operationalisiert ist.
  \item \textbf{Subagenten-L{\"u}cken:} Task~B (Simulation) war im ersten Lauf l{\"u}ckenhaft; drei fokussierte Mini-Worker~\code{B1--B3} wurden nachgezogen und lieferten zentrale Befunde.
\end{itemize}

\section{Vergleich mit dem Stand der Forschung}\label{sec:vergleich-forschung}

Die in Kapitel~\ref{ch:related} identifizierten Forschungsl{\"u}cken treffen sich mit den \agora{}-Befunden an mehreren Stellen:

\begin{itemize}
  \item \textbf{Persona-Validierung:} \agora{} implementiert DACH-Quoten-Kalibrierung (Konstanten, kein Mikrodaten-Fit) -- die Forschungsl{\"u}cke bleibt offen~\cite{uas2025,analyticflexibility}.
  \item \textbf{Provenance-Verifikation:} \agora{} erzwingt serverseitige Anker; die Verifikation {\"u}ber den Anker scheitert jedoch (F1, F3) -- direkt im Anschluss an~\cite{citednotverified}.
  \item \textbf{Multi-Agent-Debate + Quellenverifikation:} \agora{} kombiniert Red-Team und Evidence-Gating; die isolierte Wirkung der Kombination ist nicht gemessen -- Anschluss an~\cite{mast2025,hiddenbench2025}.
  \item \textbf{Model-Collapse / Drift:} \agora{} hat keine L{\"a}ngszeit-Stabilit{\"a}tsmessung -- direkte Forschungsl{\"u}cke~\cite{shumailov2024curse}.
\end{itemize}

\section{Konsequenzen f{\"u}r andere Multi-Agent-Systeme}\label{sec:konsequenzen}

Die Befunde sind nicht \agora{}-spezifisch. Sie verallgemeinern auf jedes LLM-basierte System, das eine Evidence- oder Confidence-Pipeline {\"u}ber Multi-Agent-Simulation implementiert:

\begin{enumerate}
  \item \emph{Strikte Downstream-Gates ohne Provenance-Upstream erzeugen null validierte Ausgabe -- nicht {\"u}berpr{\"u}fbare Strenge.}
  \item \emph{Echo-Caps sind nur dann wirksam, wenn sie cross-stakeholder-Konsens adressieren; intra-group-Echo-Kammern bleiben unbeobachtet.}
  \item \emph{Confidence-Konzepte m{\"u}ssen getrennt werden: synthetischer Konsens, Evidence-gest{\"u}tzte Confidence, empirische Confidence.}
  \item \emph{Hartkodierte Schwellen sind Engineering-Defaults, keine validierten Hyperparameter; ohne Kalibrierung sind sie Selbsttests.}
\end{enumerate}

\section{Methodische Lehre}\label{sec:methodische-lehre}

Der zur{\"u}ckgezogene Befund~C56~\cite{zenodo-agora} illustriert eine methodische Falle: eine einzelne Beobachtung (fehlendes Feld) wird ohne Pr{\"u}fung der Schreibpfade zur Ursachenaussage erhoben. Die Schlussfolgerung \enquote{\code{supports\_claim} im Interview-Zweig setzen} w{\"u}rde ADR-0011-Anker~3 wieder hergestellt und ADR-0002-Hartanker~5 unterlaufen haben. Die Korrektur wurde durch einen \code{grep} {\"u}ber den Schreibpfad erreicht -- und ist selbst ein Beispiel daf{\"u}r, dass Ein-Feld-Inspektion ohne Pfad-Kontext irref{\"u}hren kann.

% =====================================================================
\part{Schluss}
% =====================================================================

% =====================================================================
% Kapitel 7: Schluss
% =====================================================================

\chapter{Schluss}\label{ch:schluss}

\section{Beantwortung der Forschungsfrage}\label{sec:antwort}

Die Forschungsfrage lautete: \emph{Welche Teile von \agora{} erzeugen einen tats{\"a}chlich {\"u}berpr{\"u}fbaren Erkenntnisgewinn gegen{\"u}ber einem einzelnen starken LLM-Prompt -- und welche Teile erzeugen lediglich zus{\"a}tzliche Komplexit{\"a}t, Kosten oder scheinbare Evidenz?}

\textbf{Antwort:} Konzeptionell {\"u}berpr{\"u}fbar sind die Strukturen, die nicht prompt-basierbar sind -- Echo-Cap, zweistufiges Binding, Determinismus vor Embedding, \code{source\_kind}-Trennung, kanonische \code{evidence\_id}, Downgrade als Identit{\"a}tswechsel und die prozessualen Hartanker~\cite{adr-0002}. In der aktuellen Umsetzung (Stand~\code{7e42ae34}) ist \emph{keiner} dieser Mechanismen empirisch wirksam: der Referenzlauf~\cite{audit-r2} erzeugt null validierte Claims, weil die Provenance end-to-end gebrochen ist und der Anker opak akzeptiert wird. Damit erzeugt das System aktuell \emph{keinen} messbaren Mehrwert gegen{\"u}ber einem starken Single-Prompt -- bei gleichzeitig h{\"o}herer Komplexit{\"a}t und h{\"o}heren Kosten.

Die Antwort ist also nicht \enquote{ja, der Mehrwert ist real}, sondern \enquote{der Mehrwert ist \emph{potenziell} real, aber \emph{aktuell} nicht nachweisbar}. Genau diese Asymmetrie -- Konzept vs.\ Wirksamkeit -- ist der eigentliche Befund.

\section{Ausblick: Priorisierte Roadmap}\label{sec:roadmap}

Die priorisierte Roadmap~\cite{zenodo-agora} gliedert sich in vier Stufen:

\begin{description}
  \item[P0 -- Blocker f{\"u}r 1.0] \mbox{}\\
    \begin{itemize}
      \item ADR-0013 Teil~B umsetzen (document\_id/chunk\_id durch Graph bis Anker). Behebt F1, F3, F5, F6.
      \item Confidence-Konzepte trennen (simulation\_consensus, evidence\_confidence, empirical\_confidence). Behebt F8.
    \end{itemize}
  \item[P1 -- Empirischer Mehrwert] \mbox{}\\
    \begin{itemize}
      \item E2E-Re-Run gegen~\code{d7d9f0a4} zur Wirkungsbest{\"a}tigung des Interview-Binding-Fix. Behebt F2, G9.
      \item Baseline-Runner \agora{}-vs-Single-Prompt. Behebt G3.
      \item Sim-Validity-Metrik gegen Ground-Truth. Behebt G1.
      \item Reproduzierbarkeit (seed/determinism). Behebt G2.
    \end{itemize}
  \item[P2 -- Qualit{\"a}t] \mbox{}\\
    \begin{itemize}
      \item Markdown-Export mit Producer-Aufl{\"o}sung. Behebt F12.
      \item Hardstop bei Evidence-Omission. Behebt F13.
      \item Chunk-Provenance End-to-End-Test. Behebt G4.
      \item Negativer E2E-Test (doc\_id fallen lassen $\to$ ablehnen). Behebt G8.
    \end{itemize}
  \item[P3 -- Komplexit{\"a}ts{\"o}konomie] \mbox{}\\
    \begin{itemize}
      \item Neo4j-Evaluation (SQLite + pgvector f{\"u}r Single-User). Behebt F14.
      \item CAMEL http/oasis-Konsolidierung. Behebt F15.
      \item Redis optional markieren.
      \item Usage-Metrik {\"u}ber \code{EvidenceSourceKind}-H{\"a}ufigkeit. Behebt G7.
    \end{itemize}
\end{description}

\section{Verf{\"u}gbarkeit und Reproduzierbarkeit}\label{sec:verfuegbarkeit}

\begin{itemize}
  \item \textbf{Repository:} \href{https://github.com/arn0ld87/agora}{\code{github.com/arn0ld87/agora}}, Branch \code{feat/1152-provenance-retrieval}, HEAD \code{7e42ae34}.
  \item \textbf{Eingefrorener Snapshot:} \href{https://zenodo.org/records/21855023}{Zenodo DOI 10.5281/zenodo.21855023}~\cite{zenodo-agora}.
  \item \textbf{Autor:} Alexander Schneider, ORCID \href{https://orcid.org/0009-0008-7552-1819}{0009-0008-7552-1819}.
  \item \textbf{Audit-Dokumentation:} \code{docs/paper/agora-evidence-chain-audit.md} (43~KB), Research-Notes unter \code{docs/paper/research-notes/}, Citation-Registry~\cite{citation-registry}.
\end{itemize}

\section{Schlussbemerkung}\label{sec:schlussbemerkung}

\agora{} ist eine sorgf{\"a}ltig gebaute Maschine, die beweist, dass sie nichts beweisen kann. Das ist, paradoxerweise, genau die Ehrlichkeit, die dem Feld der LLM-basierten Sozialwissenschaft fehlt~\cite{bail2024genai,analyticflexibility}. Der Bauplan f{\"u}r den {\"u}berf{\"u}hrbaren Mehrwert liegt vor -- in~\cite{adr-0011},~\cite{adr-0013} und der Roadmap dieser Arbeit. Was fehlt, ist die Umsetzung.

% =====================================================================
\appendix
% =====================================================================

% =====================================================================
% Anhang A: Verifikationsprotokoll und Citation-Registry
% =====================================================================

\chapter{Anhang~A: Verifikationsprotokoll}\label{ch:anhang-a}

\section{Verifikations-Spot-Checks}\label{sec:verif}

Die folgenden Spot-Checks wurden exemplarisch durchgef{\"u}hrt; die vollst{\"a}ndige Liste ist in der Research-Note \code{lead-verified-findings}~\cite{zenodo-agora} enthalten.

\begin{table}[h]
  \centering
  \small
  \begin{tabular}{@{}p{0.20\textwidth}p{0.50\textwidth}p{0.20\textwidth}@{}}
    \toprule
    \textbf{Beleg} & \textbf{Inhalt} & \textbf{Verifikation} \\
    \midrule
    \src{C26}{apply\_echo\_cap} & Cap bei \code{echo\_index > 0.75} $\to$ max 0.84 / medium & \code{confidence\_calculator.py:363--386} \\
    \src{C11}{split\_text ohne Manifest} & \code{List[str]} ohne \code{doc\_id/chunk\_id} & \code{graph\_build.py:439} \\
    \src{C9}{opaker Anker} & \code{seed\_doc:}-Pr{\"a}fix akzeptiert, kein Lookup & \code{evidence.py:352} \\
    \cite{audit-r2} & 0~validierte Claims & \code{docs/reference-runs/2026-08-09-...} \\
    \cite{adr-0013} & Teil~B fehlt & ADR-0013 §1, kein \code{document\_id} in \code{graph\_build.py} \\
    \src{C45}{echo\_chamber\_index} & \code{intra/total} & \code{network\_analytics.py:426--432} \\
    \src{C46}{nicht deterministisch} & \code{random.*} ohne \code{random.seed()} & \code{run\_parallel\_simulation.py:1423,1434,1437} \\
    \src{C50}{Interview-Fix in HEAD} & \code{producer\_key} + \code{quote} + \code{persona\_stakeholder\_group} & \code{agent.py:365--413}, Test \code{test\_report\_tool\_evidence.py:110} \\
    \src{C52}{DACH-Quoten} & Destatis WZ 2008 / Mikrozensus 2024 / BFS / Statistik Austria & \code{persona\_quota\_defaults.py:1--9}, \code{persona\_demographics.py:1--6} \\
    \bottomrule
  \end{tabular}
  \caption{Verifikations-Spot-Checks (Auszug). Die vollst{\"a}ndige Liste enth{\"a}lt 33~Belege; jede Kernaussage ruht auf mindestens zwei Belegen.}
  \label{tab:verif}
\end{table}

\section{Failure-Mode-Tabelle (vollst{\"a}ndig)}\label{sec:fm-voll}

\begin{table}[h]
  \centering
  \small
  \begin{tabular}{@{}p{0.04\textwidth}p{0.40\textwidth}p{0.10\textwidth}p{0.36\textwidth}@{}}
    \toprule
    \textbf{\#} & \textbf{Failure-Mode} & \textbf{Severity} & \textbf{Beleg / Gegenmaßnahme} \\
    \midrule
    F1 & Seed-Provenance end-to-end gebrochen & Critical & \src{C11}{}; ADR-0013 Teil~B \\
    F2 & Interview$\to$Evidence-Binding-Defekt & Medium & \src{C50}{}, \code{d7d9f0a4}; E2E-Re-Run \\
    F3 & Opaker \code{seed\_doc:}-Anker & Critical & \src{C9}{}, \src{C10}{} \\
    F4 & Gate zu streng $\to$ 0~validierte Claims & High & Folge von F1+F2 \\
    F5 & Graph-Fakten ohne Dokumentidentit{\"a}t & High & \src{C19}{}, \src{C20}{}, \src{C21}{} \\
    F6 & \code{medium}-Stufe unerreichbar & High & \src{C8}{}; Folge von F1 \\
    F7 & Historische Alt-Labels: 25~medium-Claims auf 100\,\% \code{inferred} & High & \cite{adr-0011}, \cite{audit-r1} \\
    F8 & Confidence-Konzeptvermischung & Medium & \cite{audit-r2} \\
    F9 & Keine Sim-Validity-Metrik & Medium & \src{C41}{}, \src{C46}{} \\
    F10 & Keine Reproduzierbarkeit & Medium & \src{C42}{}, \src{C46}{}, \src{C47}{} \\
    F11 & Kein Baseline-Vergleich \agora{}-vs-Single-Prompt & Medium & G3 \\
    F12 & Markdown-Export ohne Producer-Aufl{\"o}sung & Low & \src{C39}{} \\
    F13 & Evidence-Map f{\"a}llt stumm bei \code{contract\_violation} & Medium & \cite{issue-987}, \src{C37}{} \\
    F14 & Neo4j overkill & Low & \src{C40}{} \\
    F15 & CAMEL http/oasis-Divergenz & Low & \src{C40}{} \\
    F16 & \code{simulated\_hours} nicht exponiert & Low & \cite{issue-1018} \\
    \bottomrule
  \end{tabular}
  \caption{Vollst{\"a}ndige Failure-Mode-Tabelle (16~Eintr{\"a}ge).}
  \label{tab:fm-voll}
\end{table}

\section{Citation-Registry}\label{sec:registry}

Die vollst{\"a}ndige Citation-Registry umfasst:

\begin{itemize}
  \item 57~Code-Belege \code{C1--C57},
  \item 5~ADRs \code{A1--A5},
  \item 12~Issues \code{I1--I12},
  \item 4~Referenzl{\"a}ufe \code{R1--R4},
  \item 14~Literaturschl{\"u}ssel \code{L1--L14}.
\end{itemize}

Jeder Beleg ist in der Research-Note \code{citation-registry.md}~\cite{citation-registry} mit Symbol, Datei, Zeilennummer und Vertrauensgrad hinterlegt. \code{README.md}, \code{docs/architecture.md} und \code{VISION.md} werden ausdr{\"u}cklich \emph{nicht} als Belege verwendet.

\section{Der zur{\"u}ckgezogene Befund C56}\label{sec:c56}

Der urspr{\"u}ngliche Befund \enquote{\code{supports\_claim} wird im Interview-Item nicht gesetzt} beruhte auf einer Fehldeutung der zweistufigen Binding-Architektur.

\begin{itemize}
  \item \code{supports\_claim} wird an genau einer Stelle gesetzt: \code{evidence\_binder.py:123} -- \code{bound["supports\_claim"] = result.verdict is EntailmentVerdict.SUPPORTED}.
  \item \code{evidence\_entailment.py:11} h{\"a}lt fest: \enquote{Nur hier darf \code{supports\_claim = True}}. \code{agent.py:649} verweist im Kommentar ausdr{\"u}cklich auf die Entailment-Stufe als setzende Instanz.
  \item Sachlich zwingend: Ein Evidence-Item ohne Claim-Bezug kann nicht beantworten, ob es \emph{einen} Claim st{\"u}tzt. Die Frage entsteht erst beim Binding.
\end{itemize}

\textbf{Schlussfolgerung:} Die Beobachtung war richtig, die Schlussfolgerung nicht. Die vorgeschlagene Ma{\"s}nahme w{\"a}re eine Regression gewesen -- sie h{\"a}tte Defekt~3 aus~\cite{adr-0011} wieder hergestellt. C56 ist damit als Beleg f{\"u}r einen \code{supports\_claim}-Defekt ung{\"u}ltig; die Methodik der Pr{\"u}fung an einer einzelnen Fundstelle wird in Kapitel~\ref{ch:diskussion} (Abschnitt~\ref{sec:methodische-lehre}) reflektiert.

\textbf{Nicht betroffen:} C56 belegt in~G6, dass \code{persona\_stakeholder\_group} eine freie Rollenbezeichnung ist und keine kodierte Taxonomie. Diese Beobachtung bleibt g{\"u}ltig und ist eigenst{\"a}ndig relevant: \code{cross\_stakeholder\_for\_high} z{\"a}hlt distinkte Werte, sodass zwei Schreibweisen derselben Rolle als zwei Stakeholder-Gruppen z{\"a}hlen.


% --- Bibliographie ----------------------------------------------------
\printbibliography[heading=bibintoc,title={Literaturverzeichnis}]

\end{document}
